
\chapter{Идентификация клиентов и~ТСП}\label{ch:kyc-kyb}
Идентификация физлица (KYC, Know Your Customer) и идентификация
юрлица или ТСП (торгово-сервисное предприятие; KYB, Know Your Business)
формируют профиль клиента (совокупность собранных и~оценённых сведений,
на которой строится решение по операциям),
от которого зависят все последующие операции.
Комплаенс в~платёжном бизнесе спрашивает: кого
подключаем, что уже известно о~клиенте или ТСП и~в~какой момент
операцию нужно остановить. Это рабочий механизм внутри операционного потока.

\section{Идентификация при подключении физлиц}\label{sec:kyc-individuals}

\highlight{Идентификация физлица -- это лестница уровней
доверия}. На каждой ступени оператор решает, можно ли считать
человека тем, за кого он себя выдаёт, и~какой риск принимает,
открывая следующий уровень сервиса. В~платёжном приложении после
регистрации часть лимитов работает сразу, а~другие открываются
после дополнительной проверки.

Лестница идёт от минимального набора сведений к~документальной
сверке и дополнительным проверкам там, где этого требует закон или
риск-профиль~--- набор признаков клиента, по которым оператор
относит его к~уровню риска, заданному Положением Банка России
№~499-П~\cite{cbr-499p}. Чем меньше сведений собрано на входе,
тем уже набор доступных операций и~лимитов; Банк России прямо
разъясняет это на примере переводов и электронных кошельков по
упрощённой идентификации~\cite{cbr-banking-faq-identification}.
Автоматизация ускоряет вход, но не отменяет обязанности проверять
качество данных: ошибка на входе редко ограничивается одним клиентом.

Стандартный набор проверок -- сверка документов
по государственным источникам, проверка по перечням
Росфинмониторинга и иным санкционным спискам (см.\ главу~\ref{ch:sanctions}),
а~также выявление политически значимого лица (ПЭП; калька
с~английского PEP, politically exposed person)~--- человека, наделённого
или ранее наделённого значимой публичной функцией; в~процедурах
противодействия легализации (отмыванию) доходов и~финансированию
терроризма (ПОД/ФТ) отдельно учитывают также членов его семьи
и~близких соратников и~применяют к~этой категории усиленную проверку
из-за повышенного коррупционного риска~\cite{cbr-499p,fatf-pep-guidance}.
Сама методология ПЭП восходит к~рекомендациям FATF; о~приостановке
членства РФ в~FATF и~текущей роли ЕАГ как региональной FATF-style body
для СНГ~--- см.~\ref{sec:aml-fatf-eag}.

Для неидентифицированного пользователя электронных денежных
средств (ЭДС) Федеральный закон №~161-ФЗ \enquote{О~национальной
платёжной системе} (статья~10) ограничивает остаток ЭДС суммой
15\,000~руб. в~любой момент и~месячный оборот -- суммой
40\,000~руб., и~эти ограничения действуют до прохождения
упрощённой идентификации (без подтверждения личности по документу
с~фотографией). После упрощённой идентификации
лимит остатка ЭДС поднимается до 100\,000~руб., а~месячный
оборот -- до 200\,000~руб.~\cite{fz-161}. Для полной
(персонифицированной, то~есть с~очной или эквивалентной сверкой
документа, удостоверяющего личность) идентификации физлица-резидента~РФ
лимит остатка установлен законом на уровне 600\,000~руб.;
дополнительные параметры могут определяться договором
с~оператором. Эти пороги задают архитектуру уровней, и~продукт
должен хранить текущий статус идентификации клиента и применять
соответствующую систему лимитов до каждой операции.

Чем выше риск,
тем меньше у~системы оснований полагаться на самодекларацию и тем
нужнее внешние подтверждения. У~ТСП задача сложнее ещё на один
шаг -- там проверяют и человека, и~бизнес как действующую структуру.

\section{Идентификация при подключении ТСП}\label{sec:kyb-merchants}

У~юридического лица проверяют и~личность, и~деловую
конструкцию: структуру собственности, реальный вид деятельности,
возвраты, спорные операции и~риск использования бизнеса как
транзитной оболочки. Платформа, подключающая нового продавца,
обязана понять, кто за ним стоит и~как он зарабатывает.
\highlight{Идентификация юрлица собирает регистрационные данные, владельцев
и платёжное поведение в~единый профиль риска}.

Базовый пакет документов включает выписку из реестра, данные
директора и учредителей, сведения о~бенефициарном владельце
(ultimate beneficial owner, UBO), банковские реквизиты и описание
бизнеса. Если в~структуре владения остаются нестыковки и~сайт не похож на
заявленный вид деятельности, это сразу повод для дополнительных вопросов. Проверка
строится на сопоставлении официальных данных с~тем, как бизнес
выглядит в~реальности. Формально чистая компания может
оказаться рискованной, поэтому идентификация ТСП не сводится к~папке
с~документами; после подключения его сверяют с~живым поведением
бизнеса (наблюдаемыми транзакционными паттернами). Логика
идентификации по Федеральному закону №~115-ФЗ \enquote{О~противодействии
легализации (отмыванию) доходов, полученных преступным путём,
и~финансированию терроризма}~\cite{fz-115} и~Положению Банка России
№~499-П это прямо предусматривает: \highlight{кредитная организация должна установить данные
самого клиента и~сведения о~бенефициарном владельце}~\cite{cbr-499p}.

Код категории торгового предприятия (Merchant Category Code, MCC)
сам по себе нейтрален; правила схем требуют
корректно классифицировать ТСП и передавать правильный код
категории. Для категорий, которые Visa относит к~high-integrity risk
merchants (категории повышенного риска по классификации Visa),
действуют отдельные ограничения и требования
мониторинга~\cite{visa-core-rules}. Подмена MCC -- попытка спрятать
реальный профиль бизнеса под более мягкой категорией. Если
заявленная категория не сходится с~фактической деятельностью,
платёжный поток останавливают до выяснения.

Модель платёжного фасилитатора (payment facilitator, PayFac)
меняет место ответственности за идентификацию суб-мерчантов
(Sponsored Merchants~--- продавцов, подключённых через PayFac),
их KYB, текущий мониторинг и~разбор споров, но не смысл проверки.
Visa возлагает
на эквайера ответственность за Payment Facilitator и его Sponsored
Merchants: до начала обработки нужно подтвердить регистрацию PayFac
у Visa, присвоить уникальные идентификаторы самому PayFac и~его
Sponsored Merchants, получать отчётность по их активности и следить
за корректным MCC~\cite{visa-core-rules}. Платформа подключает
субмерчантов быстрее, но расплачивается собственными процедурами идентификации
и~более плотным мониторингом. Без них быстрое подключение
превращается в~быстрое накопление риска.

\practicenote{%
  При проверке ТСП ключевой сигнал~--- несоответствие заявленного MCC
  фактической деятельности.
}
\begin{table}[H]
  \begingroup
  \scriptsize
  \setlength{\tabcolsep}{4pt}
  \renewcommand{\arraystretch}{1.2}
  \begin{minipage}{\textwidth}
    \begin{tabularx}{\textwidth}{
        @{}B{0.18\textwidth}
        Y
        Y
        Y@{}
      }
      \toprule
      \headrow
      \thd{Критерий} &
      \thd{Низкий риск} &
      \thd{Средний риск} &
      \thd{Высокий риск} \\
      \midrule
      Примеры MCC &
      retail, SaaS (Software as a Service) &
      travel &
      crypto exchange,\newline gambling, adult \\
      Документы &
      стандартный пакет &
      расширенный пакет\newline + лицензии &
      расширенный пакет\newline + финансовая отчётность \\
      Мониторинг после подключения &
      квартальный &
      ежемесячный &
      еженедельный \\
      Плавающий резерв &
      нет &
      5--10\% &
      10--20\% \\
      \bottomrule
    \end{tabularx}
    \par\smallskip
    \tablenote{Цветом показана глубина должной проверки
      (due diligence)~--- комплекса процедур по~сбору и~оценке
      сведений о~клиенте: чем выше риск-профиль ТСП, тем шире пакет
      документов, плотнее мониторинг после подключения и~вероятнее
      плавающий резерв (\textit{rolling reserve})~--- удержание
      части оборота ТСП на~отдельном счёте до~закрытия диспутного
      окна.}
  \end{minipage}%
  \endgroup
  \caption{Список для идентификации ТСП по категориям риска.}\label{fig:kyb-risk}
\end{table}

\section{Мониторинг ТСП после подключения}\label{sec:merchant-monitoring}

Риск меняется после подключения, поэтому мониторинг не заканчивается
в~момент подключения. Схемы типа bust-out (когда клиент или~ТСП длительное время ведёт
себя добросовестно, чтобы нарастить лимиты, а~затем резко уводит
средства; подробнее~--- в~гл.~\ref{ch:fraud}) и~попытки использовать
ТСП для~отмывания
проявляются именно здесь. Пока профиль выглядит спокойно, мониторинг
ослабляет внимание, а~проблема становится видна поздно. Мониторинг
отслеживает изменения объёма, географии и доли возвратов в~реальном
времени.

Поведенческие признаки для расследования~---
резкий рост оборота, нетипичная география клиентов, высокий процент
возвратов и~несоответствие MCC реальной деятельности. Без контекста даже чистая статистика вводит
в~заблуждение.

В~промышленных системах это надстройка над авторизационной базой:
агрегаты по ТСП, скользящие окна, поиск выбросов и~очередь случаев
для аналитиков. Резкий рост объёма у~ТСП команда видит почти сразу,
а~не в~конце месяца. В~европейском регионе Visa требует от эквайера
сохранять ежедневные данные по обороту, среднему чеку, количеству
операций, времени до расчёта и количеству диспутов и сравнивать их
с~нормальным профилем активности минимум ежедневно~\cite{visa-core-rules}.
Чем больше признаков, тем важнее отсекать шум.
